\documentclass[12pt]{article}
\usepackage[utf8]{inputenc}
\usepackage[T1]{fontenc}
\usepackage{hyperref}
\title{Sawi Language Resources for K-12 Curriculum Development}
\author{Compiled by Claude Sonnet 5.5 (claude-sonnet-5-5) for Andr\'as Kornai}
\date{}
\begin{document}
\maketitle

\section{Introduction}
Sawi (Sawuy; ISO 639-3 saw) is a Trans-New Guinea language of the southern lowlands of Papua, Indonesia, closest to the Awyu languages. Estimates are about 3,500 to 4,000 speakers (a 1993 figure) across several villages in Asmat Regency. It should not be confused with Siawi of Papua New Guinea. Little published material was found beyond Bible translation.

\section{Primary Educational Resources}
No primers, textbooks or school programmes were found.

\section{Linguistic and Reference Materials}
\begin{itemize}
\item Wikipedia on the Sawi language: \url{https://en.wikipedia.org/wiki/Sawi_language_(Papuan)}
\item Sawi numerals page, MPI EVA: \url{https://lingweb.eva.mpg.de/channumerals/Sawi.htm}
\item University of Zurich New Guinea World page: \url{https://newguineaworld.linguistik.uzh.ch/families/trans-new-guinea/central-west-new-guinea/digul-river-ok/digul-river/sawi}
\item A vocabulary and religious-terms study (Yost, 1988) and SIL survey word lists were mentioned in secondary sources. No link was seen and they are unconfirmed.
\end{itemize}

\section{Religious and Community Resources}
The New Testament was reportedly completed in 1994, with the Old Testament in progress. Wikipedia credits Don Richardson with starting the translation, but I could not confirm this from a second source. A Mission Network News report describes churches gathering for the Old Testament project.
\begin{itemize}
\item \url{https://www.mnnonline.org/?p=170181}
\item Wikipedia on the Sawi people: \url{https://en.wikipedia.org/wiki/Sawi_people}
\item Joshua Project profile (Sawuy): \url{https://www.joshuaproject.net/index.php/people_groups/18881}
\end{itemize}

\section{Digital and Audio Resources}
None confirmed. The New Testament is reported to be on YouVersion, but no URL was seen.

\section{Curriculum Development Context}
Sawi has few speakers, and older adults are reported to be its main users. No school use is known, and no dictionary was found. The Bible translation is the main written corpus and could be a source of text, but any curriculum would need new documentation and materials work with the community.
\end{document}